\documentclass[11pt]{article}
\usepackage[margin=0.8in]{geometry}
\usepackage[]{xcolor}
\usepackage{amsmath}
\parindent=0pt

\title{File Systems}
\author{gordonbchen}
\date{}

\begin{document}
\maketitle

\section{Memory}
Dynamic memory (DRAM): capactior charged/discharged

Magnetic tape / disk: flux direction north/south

Read-only compact disct: no pit, pit, laser burns pit, reflection angle changes


Volatile: loses data on turn off


4x32 RAM: 4 addressable locations, 32 bits at each location

SRAM: flip-flops

DRAM: capacitors, must be refreshed

SSD: flash memory, P(rogram) / E(rase) cycle limit

Magnetic HDD: magnetic spinning thing

\section{File systems}
\begin{itemize}
    \item Sector 0 Master Boot Record, partition table, partitions
    \item Partition: boot block, superblock, free space mgmt, inodes, ...

    \item BIOS executes MBR program. MBR program finds active partition (where OS is), runs boot block,
        boot block loads OS into RAM.

    \item Contiguous allocation: removing files creates holes, can't grow files

        Good for CD-Rom b/c never change memory

    \item Linked-list allocation: allocate in blocks

        files are scattered through the system: pointer to physical block and next physical block

        sequential (not random) access b/c have to go through block pointers

        seeing which block is hard b/c block size is not a power of 2 b/c of pointer

    \item FAT: file allocation table (windows)

        start of partition

        in directory structure: store 1st block num. then fat stores next block, table[block idx] = next idx.

        FAT[last block] = -1

        requires huge FAT: 2GB disk, 1kb block = 200 M entries

        fixed size based on partition size and block size

    \item I-nodes (linux)

        file attributes

        first 10 addresses of data blocks (based on avg file size)

        address of disk block containing additional pointers to indirect blocks

        indirect blocks points to data blocks or other pointer blocks

        1 single indirect block + 1 double indirect block + 1 triple indirect block \\

        ex: 10 initial addr, 1kb blocks, 4-byte pointer to data block

            max file size = $10*1Kb + (1024/4)*1Kb + (256^2)*1Kb + (256^3)*1Kb$

\end{itemize}
\end{document}
